\section{Sphericalization, central induction, and rigidity}\label{sec:core}
Throughout, a fusion category is finite, semisimple, rigid and
$\CC$-linear, with simple tensor unit. Its rank is the number
of isomorphism classes of simple objects. A fusion ring is a based
ring with nonnegative integral structure constants and the usual
coefficient-of-unit pairing. A \emph{Grothendieck ring} is a fusion
ring realized by such a category. A pivotal structure is a monoidal isomorphism from the identity to
the double dual; it is spherical when its left and right traces agree.
We impose no pivotal hypothesis; the general existence question is
\cite[Conjecture~2.8]{ENO}.

For a self-dual basis $1,X,Y,Z$ we use the convention
\begin{equation}\label{eq:tuple-convention}
\begin{aligned}
X^2&=1+aX+bY+cZ,&XY&=bX+dY+eZ,&XZ&=cX+eY+fZ,\\
Y^2&=1+dX+gY+hZ,&YZ&=eX+hY+iZ,&Z^2&=1+fX+iY+jZ.
\end{aligned}
\end{equation}
Thus a table is denoted by $(a,b,c,d,e,f,g,h,i,j)$.
The same convention applies to signed tables. All subsequent changes
of labels or signs are explicitly local to the argument concerned.

\begin{lemma}[Rank-four algebras and coefficient bounds]
\label{core:lem:rank-four-algebra}\label{core:lem:coefficient-bound}
A semisimple complex algebra of dimension four admitting a
one-dimensional representation is commutative. In particular every
rank-four fusion ring, and every rank-four signed sphericalization
quotient, has commutative complexification. For an ordinary fusion ring,
write $d_i=\FPdim(b_i)$ and $D=\sum_jd_j^2$. Then
\[
 N_{ij}^{k}\le d_i<\sqrt D.
\]
\end{lemma}
\begin{proof}
The only noncommutative semisimple complex algebra of dimension four
is $M_2(\CC)$, which has no one-dimensional representation.
The regular based $*$-algebras in question are semisimple; their
Frobenius--Perron or signed dimension characters supply the required
representation. In the ordinary ring, reciprocity gives
$N_i^*=N_{i^*}$, so commutativity makes $N_i$ normal. Its operator norm
is its spectral radius $d_i$. Every matrix entry is bounded by that
norm, and $d_i^2<D=\sum_jd_j^2$.
\end{proof}

For a commutative self-dual table its characters have real values,
and its formal codegrees are
$f_\chi=\sum_i\chi(X_i)^2$. Orthogonality gives
\begin{equation}\label{core:orthog}
\sum_\chi\frac{\chi(X_i)\chi(X_j)}{f_\chi}=\delta_{ij},
\qquad\sum_\chi f_\chi^{-1}=1.
\end{equation}
All codegrees are positive. The $d$-number and cyclotomicity conditions
are collected in Proposition~\ref{finite:prop:standard}. In degree two,
the $d$-number criterion says that the norm divides the square of the
trace. A real character field of a categorification is totally real
and abelian.

We write $R=K_0(\C)$, $d_i=\FPdim(X_i)$ and
$D=\sum_i d_i^2$. In the self-dual case, let $\mathcal S$ denote
the sphericalization and $\varepsilon$ its canonical central sign
object. Its central lift is the nontrivial object of the canonical
$\Rep(C_2)$ in the equivariantization construction of sphericalization
\cite[Section~5.3]{Bartlett}; its spherical dimension and ribbon twist
are both $-1$. Write $T=K_0(\mathcal S)/(\varepsilon+1)$ and
$x_i=\dim(X_i^+)$, so $E=\sum_i x_i^2$ and
$\dim\mathcal S=2E$. The lifts are self-dual: duality preserves
spherical dimensions, and the two lifts have opposite nonzero dimensions.
The notation $\beta_i=x_i$ is used when the lift signs are chosen
to make the embedded spherical dimensions positive at the selected real
embedding.
A Galois change of embedding is made explicitly; neither sphericality
nor a unitary realization alone is used to assert pseudo-unitarity. The ordinary and signed
codegrees are denoted by $f_i,g_i$, and
\[
u=\sum_i f_i^{-2},\qquad v=\sum_i g_i^{-2}.
\]
\begin{lemma}[Sphericalization and induction]\label{core:lem:family-induction}
Let $\C$ be a complex fusion category with a self-dual
rank-four ring $R$, FP dimensions $d_i$, and $D=\sum_i d_i^2$.
Its sphericalization has a commutative signed quotient $T$ of rank
four. Choose positive embedded spherical dimensions $\beta_i$ for its lifts and
put $E=\sum_i\beta_i^2$. Then:
\begin{enumerate}
\item $|T_{ij}^k|\le N_{ij}^k$ and $T_{ij}^k\equiv N_{ij}^k\pmod2$;
\item all ordinary and signed codegrees lie in $L=\Q(\beta_i)$,
and $E/f$ is an algebraic integer for every such codegree $f$;
\item every signed codegree is at most $D$ and $E\le D$;
\item if $E=D$, the chosen lifts give a spherical section of
$\C$.
\end{enumerate}
These statements remain valid after real embeddings of $L$ and
changes of lift signs. In particular, in seeking a nonspherical
categorification one may assume $E<D$ and select the larger
conjugate when $[L:\Q]=2$.
\end{lemma}
\begin{proof}
The sphericalization $\mathcal S$ has two self-dual lifts of each
simple, exchanged by its canonical central sign object $\varepsilon$,
whose dimension is $-1$ and whose square is one; see
\cite[Section~5.3]{Bartlett} and \cite[Section~5.1]{Ostrik3}.
Taking the quotient by $\varepsilon+1$ gives signed coefficients
$T_{ij}^k=N_{ij}^{k,+}-N_{ij}^{k,-}$, while
$N_{ij}^k=N_{ij}^{k,+}+N_{ij}^{k,-}$. Its real regular matrices
are symmetric. Semisimplicity follows from the positive
coefficient-of-unit form; the dimension character excludes
$M_2(\CC)$, the sole noncommutative semisimple algebra of
dimension four. Thus $T$ and $K_0(\mathcal S)$ are commutative.
The latter has codegrees $2f_j,2g_j$ and spherical dimension $2E$.

Central induction into $Z(\mathcal S)$ has one simple constituent
of $I(1)$ per character, of dimension $E/f_j$ or $E/g_j$ and twist
one \cite[Theorems~2.5 and~2.13]{Ostrik3}. These dimensions are
algebraic integers; codegree-field containment is
\cite[Corollary~2.15]{Ostrik3}. The signed $g_j$ are $d$-numbers
because $2g_j$ are codegrees of $\mathcal S$.
Spectral comparison $|T_i|\le N_i$ bounds every signed character
value in absolute value by $d_i$. Thus its codegree is at most $D$.
If $E=D$, then $\beta_i=d_i$ throughout, and
\[
 \sum_k(N_{ij}^k-T_{ij}^k)d_k=0
\]
forces $T_{ij}^k=N_{ij}^k$. The positive lifts are tensor closed
and form the required spherical section. The same algebraic
arguments hold after each real embedding, followed by choosing
the signs of the lifts to make their values positive.
\end{proof}

The central quotients $E/f_i,E/g_i$ belong to the signed order itself,
not merely to its field's ring of integers: the forgetful image is an
integral signed combination of the chosen lifts.
For a field algebra, the coefficient-of-unit functional satisfies
\begin{equation}\label{core:frobenius-trace}
\tau_T(a)=\Tr_{L/\Q}(a/E).
\end{equation}
Indeed $\Tr(a)=\sum_i\tau_T(x_i a x_i)=\tau_T(aE)$.
In particular $\Tr(x_i/E)=\delta_{i0}$.
When both ordinary and signed algebras are the same quartic field,
\begin{equation}\label{core:discriminants}
\disc(R)=N(D),\quad \disc(T)=N(E),\quad N(E/D)\in\ZZ_{>0}^2.
\end{equation}
The first two identities follow by comparing the trace pairing with
the orthonormal based pairing. The last follows because the ratio
of two order discriminants is a rational square and $E/D$ is integral.

\subsection{Spectral and trace saturation}


A self-dual based ring has a basis $b_0=1,b_1,b_2,b_3$ and a symmetric
integral multiplication tensor, with coefficient-of-unit form
$\langle b_i,b_j\rangle=\delta_{ij}$. A signed table below has the same
unit and form; its nonunit structure constants need not be positive.
A gauge change replaces each $b_i$ by $s_i b_i$, where $s_0=1$ and
$s_i\in\{1,-1\}$.



\begin{lemma}\label{trace:lem:absolute-rigidity}
Let $N=(N_{ijk})$ be the structure tensor of a self-dual rank-four fusion ring, and let $T=(T_{ijk})$ be a symmetric integral structure tensor defining an associative unital algebra on the same orthonormal basis $1,X,Y,Z$, with $T_{0ij}=\delta_{ij}$. If
\[
|T_{ijk}|=N_{ijk}\qquad(0\le i,j,k\le3),
\]
then there are signs $s_0=1$, $s_1,s_2,s_3\in\{1,-1\}$ such that
\[
T_{ijk}=s_i s_j s_k N_{ijk}.
\]
\end{lemma}
\begin{proof}
Associativity and symmetry give, for every $i,j$,
\[
\sum_k T_{iik}T_{jjk}=\sum_k T_{ijk}^{\,2}.
\]
The same identity holds for $N$, and the right-hand sides agree. Each term
$N_{iik}N_{jjk}-T_{iik}T_{jjk}$ is nonnegative; therefore each is zero. For a fixed $k$, all nonzero coefficients $T_{iik}$ consequently have the same sign.

Use the usual ten-entry convention
\[
(a,b,c,d,e,f,g,h,i,j)
=(N_{111},N_{112},N_{113},N_{122},N_{123},
N_{133},N_{222},N_{223},N_{233},N_{333}).
\]
The three groups of diagonal coefficients are
\[
(a,d,f),\qquad(b,g,i),\qquad(c,h,j).
\]
Independent sign changes of $X,Y,Z$ make every signed coefficient in these groups nonnegative. Hence the two tables now agree except possibly at $e$. If they agree there too, we are done. Otherwise $e>0$ and the signed table has coefficient $-e$.

Compare the associativity equations for the two choices of sign of $e$. Their linear and constant parts give
\begin{gather*}
a=d+f,\qquad g=b+i,\qquad j=c+h,\\
bh+ci=bc,\qquad dh=cd+fh,\qquad bf+di=fi,\\
e^2=1+df+bi+ch.
\end{gather*}
We claim that one of the three displayed diagonal groups is zero.
If $d>f$, then $bf=i(f-d)$ forces $i=0$ and $bf=0$.
For $f>0$ this gives $b=g=i=0$.
For $f=0$, the equation $dh=cd$ gives $h=c$; then $e^2=1+c^2$ forces $c=h=j=0$ by integrality.
If $d=f>0$, the equations give $b=c=0$ and then $e^2=1+d^2$, an impossibility.
If $d=f=0$, then $a=d=f=0$.
Finally, if $d<f$, the identity $cd=h(d-f)$ forces $h=0$ and $cd=0$.
For $d>0$ this gives $c=h=j=0$.
For $d=0$, the equation $bf=fi$ gives $b=i$; then $e^2=1+b^2$ forces $b=g=i=0$.
This proves the claim in every case.

Flip the basis sign belonging to a zero group. It preserves all nine diagonal coefficients and reverses the sign of $e$. The two tensors now agree, proving the lemma.
\end{proof}

\begin{corollary}\label{trace:cor:trace-rigidity}
Under the same hypotheses, suppose only that $|T_{ijk}|\le N_{ijk}$. Let $C_N$ and $C_T$ be multiplication by the respective Casimir elements $\sum_i b_i^2$. If
\[
\operatorname{Tr}(C_T)=\operatorname{Tr}(C_N),
\]
then $T$ and $N$ are gauge equivalent.
\end{corollary}
\begin{proof}
In the regular representation the multiplication matrices are symmetric, and
\[
\operatorname{Tr}(C_N)=\sum_{i,j,k}N_{ijk}^2,
\qquad
\operatorname{Tr}(C_T)=\sum_{i,j,k}T_{ijk}^2.
\]
Equality, together with termwise domination, forces equality of every absolute value. Apply Lemma~\ref{trace:lem:absolute-rigidity}.
\end{proof}


\begin{corollary}\label{core:trace-spherical}
Equality of the ordinary and signed Casimir traces in a rank-four
self-dual sphericalization implies that the original category is spherical.
\end{corollary}
\begin{proof}
Corollary~\ref{trace:cor:trace-rigidity} supplies a basis sign change
for which the signed and ordinary tensors agree. The chosen lifts
are tensor closed, so the spherical structure restricts to a tensor
subcategory equivalent to the original category.
\end{proof}

\begin{lemma}[Galois transport]\label{core:lem:galois-transport}
If a spherical dimension character belongs to the Galois orbit of the
Frobenius--Perron character, the category has a pseudo-unitary Galois
conjugate with the same fusion rules. Its central identities can
therefore be evaluated with positive Frobenius--Perron dimensions.
\end{lemma}
\begin{proof}
Extend the required embedding of the dimension field to an automorphism
of $\CC$. Apply it to skeletal tensor, duality and spherical
structure constants. All their identities and the fusion rules persist,
while the dimension character becomes the Frobenius--Perron character.
This requires no algebraic field of definition for the associativity
coefficients. Equivalently, apply the automorphism to the central
identities; restriction multiplicities are fixed and twists remain
roots of unity.
\end{proof}

\begin{corollary}[Multiplicity-one sphericality]
\label{core:cor:multiplicity-one-spherical}
Every complex categorification of a self-dual rank-four fusion ring
of multiplicity one admits a spherical structure.
\end{corollary}
\begin{proof}
For its signed quotient, domination and parity give
$|T_{ij}^{k}|\le N_{ij}^{k}$ and
$T_{ij}^{k}\equiv N_{ij}^{k}\pmod2$. Since $N_{ij}^{k}\in\{0,1\}$,
$|T_{ij}^{k}|=N_{ij}^{k}$ for every coefficient.
Lemma~\ref{trace:lem:absolute-rigidity} makes the two tables gauge
equivalent. The corresponding tensor-closed lifts form a spherical
subcategory equivalent to the original category.
\end{proof}

\begin{proposition}[Established inputs]\label{finite:prop:standard}
\begin{enumerate}
\item In any complex categorification every formal codegree is a
$d$-number. For its monic minimal polynomial
$t^e+a_1t^{e-1}+\cdots+a_e$, this means
$a_e^j\mid a_j^e$ for $1\leq j\leq e$
\cite{OstrikCodegrees}.
\item If $S\subseteq R$ is a fusion subring, $D/\FPdim(S)$ is an
algebraic integer \cite[Theorem 7.17.6]{EGNO}.
\item Every character value of a complex-categorifiable fusion ring is
cyclotomic \cite[Theorem 8.51 and Corollary 8.53]{ENO}.
\item In a spherical categorification with dimension character $\chi$,
every $\chi(b_i)$ is nonzero, $\chi(b_i)=\chi(b_{i^*})$, and its global
dimension $\Delta=f_\chi$ satisfies $\Delta/f\in\overline{\ZZ}$
for every formal codegree $f$. Its dimension field contains all formal
codegrees \cite[Theorem 2.13 and Corollaries 2.14--2.15]{Ostrik3}.
\item A weakly integral fusion category is pseudo-unitary
\cite[Proposition 8.24]{ENO}. A categorifiable fusion ring with precisely
two distinct simple Frobenius--Perron dimensions has a pseudo-unitary
Galois conjugate \cite[Lemma 3.1]{Schopieray}.
\end{enumerate}
\end{proposition}

For reference, for a spherical category of global dimension $H$,
central induction and higher indicators satisfy
\begin{equation}\label{core:central-indicator}
\sum_{A\in\Irr Z(\mathcal A)}[I(V):A]\dim(A)\theta_A^n=H\nu_n(V).
\end{equation}
Here the twists are roots of unity, and
$\nu_n(V)\in\Q(\zeta_n)$; see \cite[Theorem~4.1]{NS} and
\cite[Theorem~5.1]{NSHigher}. Induction of the unit has one
multiplicity-one constituent per character, of dimension $H/f_\chi$
and twist one \cite[Theorems~2.5 and~2.13]{Ostrik3}.
\begin{lemma}[The central sign and its translate twists]
\label{core:lem:translate-twists}
Let $\mathcal S$ be a rank-four self-dual sphericalization, with
central sign $\varepsilon$, and let $A_\chi\subset I(1)$ be the
multiplicity-one constituent corresponding to a character $\chi$.
Its canonical central sign translate satisfies
\[
 \theta_{\widehat\varepsilon\otimes A_\chi}=\chi(\varepsilon).
\]
Thus ordinary constituents have translate twist $+1$ and signed
constituents have translate twist $-1$. The constituents of $I(1)$
and their sign translates are disjoint. Moreover
$\nu_n(\varepsilon)=0$ for odd $n$ and $-1$ for even $n$.
\end{lemma}
\begin{proof}
Central induction identifies modular characters by
\[
 \frac{S_{A_\chi,W}}{S_{A_\chi,0}}=\chi(F(W))
\]
\cite[Example~2.9 and the paragraph preceding Theorem~2.13]{Ostrik3}.
At $W=\widehat\varepsilon$ the ratio is $\chi(\varepsilon)$;
it is also $\dim(\varepsilon)$ times the monodromy scalar.
Since $\dim(\widehat\varepsilon)=\theta_{\widehat\varepsilon}=-1$
and $\theta_{A_\chi}=1$, balancing gives the asserted twist.
The constituents have dimensions $E/f_i,E/g_i>0$, whereas their
translates have the opposite dimensions, proving disjointness.
For the indicator, the invariant space is zero for odd tensor powers
of $\varepsilon$. For even powers it is one-dimensional, and the
rotation is multiplied by the inverse pivotal map $-1$.
\end{proof}

\begin{lemma}\label{core:lem:adjoint-parity}
In a pivotal fusion category, a self-dual simple with second indicator $-1$ has even multiplicity in every adjoint product $Y^*\otimes Y$.
\end{lemma}
\begin{proof}
Let $j$ be the pivotal structure and choose $f:X\to X^*$ with transpose $f^t=f^*j_X=-f$. On $W=Y\otimes X\otimes Y^*$, the map
$\alpha=j_Y\otimes f\otimes\operatorname{id}_{Y^*}:W\to W^*$ satisfies $\alpha^t=-\alpha$, by monoidality and $j_Y^*j_{Y^*}=1$. The pairing
$B(v,w)=v^*\alpha w$ on $\operatorname{Hom}(1,W)$ is nondegenerate by semisimplicity and alternating by this transpose identity. Its dimension is even. Rigidity identifies that dimension with the desired multiplicity.
\end{proof}

